% !TeX program = xelatex
% !TeX encoding = utf8
% !TeX spellcheck = uk_UA
% !BIB program = biber
\documentclass{Syllabus}
\usepackage{needspace}
\usetikzlibrary{arrows.meta, shapes.geometric}
\addbibresource{../PlasmaBook/bibliography/bibliography.bib}

% ================================================
%                Дані  дисципліни
% ================================================
\creator{доцент, к.ф.-м.н., доцент Пономаренко Сергій Миколайович}
\lecturer[s.ponomarenko@kpi.ua]{доцент, к.ф.-м.н., доцент Пономаренко Сергій Миколайович}
\practicer[]{доцент, к.ф.-м.н., доцент Пономаренко Сергій Миколайович}
\disciplinename{Фізика плазми} % Назва дисципліни
%\disciplinecode{ПВ}
\kursurl{https://classroom.google.com/u/2/c/NjU4MTMxNzc0MTkw}
\shedule{http://ipt.kpi.ua/navchalnij-protses}
% Дати й номери протоколів — спільний файл protocols.tex
% Дати й номери протоколів — спільні для всіх силабусів кафедри.
% Підключається з силабуса: % Дати й номери протоколів — спільні для всіх силабусів кафедри.
% Підключається з силабуса: % Дати й номери протоколів — спільні для всіх силабусів кафедри.
% Підключається з силабуса: \input{protocols.tex}  (змінювати тут, не в силабусах)
% -------------- Ухвалення кафедрою ---------------
\kafday{10}
\kafmonth{червня}
\kafyear{2026}
\kafnum{7}
% -------------- Затвердження метод. комісією ---------------
\methodcomday{29}
\methodcommonth{червня}
\methodcomyear{2026}
\methodcomnum{6}
  (змінювати тут, не в силабусах)
% -------------- Ухвалення кафедрою ---------------
\kafday{10}
\kafmonth{червня}
\kafyear{2026}
\kafnum{7}
% -------------- Затвердження метод. комісією ---------------
\methodcomday{29}
\methodcommonth{червня}
\methodcomyear{2026}
\methodcomnum{6}
  (змінювати тут, не в силабусах)
% -------------- Ухвалення кафедрою ---------------
\kafday{10}
\kafmonth{червня}
\kafyear{2026}
\kafnum{7}
% -------------- Затвердження метод. комісією ---------------
\methodcomday{29}
\methodcommonth{червня}
\methodcomyear{2026}
\methodcomnum{6}

% ---------------- Syllabus head -----------------
\efield{E Природничі науки, математика та статистика}
\speciality{E6 Прикладна фізика та наноматеріали}
\kafname{прикладної фізики}
% ---------------- Syllabus head -----------------
\teachingprogramm{Прикладна фізика}
%\cycle{Цикл загальної підготовки}
\level{Другий (магістерський)} % освітньо-кваліфікаційний рівень
\status{Вибіркова}
%-------- Семестр 3 --------------------
\course{1}
\semestr{весняний}
\controltype{Залік}
\AudDistribHours%
{4}%
{30}%
{16}%
{}%


\MKR{1}
\DKR{}
\RGR{}
\Referat{}
% ================================================


\def\zsrs{\textit{Завдання на СРС:}}
\def\lit{\textit{Опрацювати:\ }}
\def\probl{\textit{Задачі:\ }}
\def\lab{\textit{Лабораторний практикум:\ }}


\begin{document}
\printhead


\syllabuschapter{Реквізити навчальної дисципліни}
\bluetableprint

\syllabuschapter{Програма навчальної дисципліни}

\section{Опис навчальної дисципліни, її мета, предмет вивчання та результати навчання}


Метою курсу є підвищення загального рівня знань здобувачів вищої освіти, введення їх у коло явищ, які розглядає фізика плазмі, ознайомлення їх з основними термінами, поняттями, методами і ідеями, що були розвинуті в цій науці, підготовка їх до самостійної роботи із оригінальною науковою літературою.

Предметом вивчення є характерні явища в системах заряджених частинок: особливості колективного руху заряджених частинок завдяки їх електромагнітній взаємодії, та особливості індивідуального руху окремих заряджених частинок в електричних та магнітних полях.


%Після засвоєння навчальної дисципліни студенти мають продемонструвати такі \emph{результати} навчання:
%
%\begin{enumerate}
%	\item[\bfseries знання:] загальних і характерних властивостей плазми як системи заряджених частинок з далекодійним характером взаємодії; основні поняття про природу елементарних та колективних процесів в плазмі, явищ переносу в плазмі за відсутності та наявності магнітного поля; особливості руху заряджених частинок в магнітному полі;  фізичні механізми електричного пробою та основні властивості плазми газового електричного розряду
%	\item[\bfseries уміння:] орієнтуватися в термінології, основах математичного формалізму модельних концепцій; виявляти і прогнозувати специфічну поведінку систем заряджених частинок, що
%	      зумовлюється далекодійною кулонівською та електромагнітною взаємодією.
%	\item[\bfseries досвід:] вільно орієнтуватися на якісному й кількісному рівні в основних фізичних явищах, пов’язаних з проявами електромагнітного поля в плазмі, виробити навички практичного використання засвоєних знань, методів і підходів у подальшому навчанні та професійній діяльності.
%\end{enumerate}

Згідно з вимогами освітньо-професійної програми здобувачі вищої освіти після засвоєння навчальної дисципліни <<\discipline>> мають продемонструвати такі результати навчання:

\subsection*{Загальні компетентності ОПП}

\begin{enumerate}
	\item [ЗК 1:] Здатність до абстрактного та аналітичного мислення, розуміння основних концепцій, парадигми та ідей прикладної фізики.
	\item [ЗК 2:] Здатність до навчання та самонавчання шляхом пошуку, аналізу та конструктивного синтезу інформації з різних джерел.
	\item [ЗК 8:] Здатність до кваліфікованого проведення досліджень на відповідному рівні під керівництвом фахівців, включаючи аналіз проблем, постановку цілей і завдань, вибір методів дослідження та аналіз отриманих результатів.
\end{enumerate}

\subsection*{Фахові компетентності ОПП}

\begin{enumerate}
	\item [ФК 2:] Здатність до безперервного поглиблення фундаментальних знань та систематичного вивчення та аналізу нової науково-технічної інформації, світового досвіду в галузі прикладної фізики та наноматеріалів.
	\item [ФК 3:] Здатність застосовувати теоретичні знання для аналізу фізичних систем, явищ і процесів в галузі прикладної фізики та наноматеріалів.
	\item [ФК 8:] Здатність використовувати методи і засоби математичного моделювання для опису фізичних об'єктів та процесів.
	\item [ФК 9:] Здатність використовувати знання про фізичну природу об'єктів у роботах по створенню нових приладів, апаратури, систем, обладнання, речовин і матеріалів (наноматеріалів).
	\item [ФК 11:] Здатність до вибору методів дослідження структури, складу та властивостей матеріалів (наноматеріалів), що використовуються або застосовуються в фізичних, біофізичних та енергетичних системах, вибору оптимальних параметрів дослідження і розуміння границь застосування обраного методу.
\end{enumerate}

\subsection*{Програмні результати навчання з ОПП}

\begin{enumerate}
	\item [ПРН 1:] Знання окремих розділів прикладної фізики на рівні, необхідному для виконання експериментальних досліджень та аналізу результатів в контексті існуючих теорій за умов невизначеності і неповноти експериментальних даних.
	\item [ПРН 2:] Знання методів теоретичної фізики, спеціальних розділів вищої математики, програмування, прикладних програм і методів обчислення на рівні, необхідному для аналізу і моделювання фізичних процесів і систем.
	\item [ПРН 4:] Знання методів теоретичної фізики, спеціальних розділів вищої математики на рівні, необхідному для розуміння функціонування та моделювання процесів, що відбуваються в технологічних та технічних системах, в тому числі інформаційних.
	\item [ПРН 9:] Вміння застосовувати фізичні, математичні та комп'ютерні моделі для дослідження фізичних явищ, розробки приладів, нових матеріалів і наукоємних технологій в області біофізики, енергетичних та інформаційних систем (залежно від освітньої траєкторії).
	\item [ПРН 10:] Вміння вибирати методи та інструментальні засоби проведення досліджень в одній із галузей прикладної фізики (залежно від освітньої траєкторії), користуватись обладнанням та устаткуванням для здійснення фізичного експерименту.
	\item [ПРН 13:] Вміння використовувати сучасні методи і технології наукової комунікації українською та іноземною мовами, вміння читати та розуміти фахові англомовні джерела.
\end{enumerate}


\section{Пререквізити та постреквізити дисципліни (місце в структурно-логічній схемі\\ навчання за відповідною освітньою програмою)}

%Зазначається перелік дисциплін, або знань та умінь, володіння якими необхідні студенту (вимоги до рівня підготовки) для успішного засвоєння дисципліни (наприклад, «базовий рівень володіння англійською мовою не нижче А2»). Вказується перелік дисциплін які базуються на результатах навчання з даної дисципліни.

Для опанування дисципліни здобувачі вищої освіти повинні володіти базовими знаннями з вищої математики (математичний і тензорний аналіз) та загальної й теоретичної фізики, зокрема розуміти основи механіки, термодинаміки й молекулярної фізики, електрики та магнетизму, оптики, теорії поля та електродинаміки суцільних середовищ.

Також повинні вміти використовувати математичний апарат: операції з матрицями,  диференціювати, інтегрувати.

Засвоєні теоретичні знання та отримані практичні навички під час вивчення навчальної дисципліни «Фізика плазми» можна використовувати в навчальних дисциплінах, пов’язаних з вивченням явищ, технологій та методик, в яких використовуються або розглядаються явища пов’язані із плазмою, або з системами заряджених частинок (плазмові та іонно/електронно променеві технології та експериментальні методики/обладнання).

Взаємозв'язок тем самої дисципліни показано на рис.~\ref{fig:topics-graph}: стрілка веде від теми, на яку спирається виклад, до теми, що її використовує. Номери тем збігаються з номерами рядків таблиці в розділі <<Методика опанування навчальної дисципліни>>.

\begin{figure}[h!]
\centering
\begin{tikzpicture}[
	x=1cm, y=1cm,
	node distance=0pt,
	t/.style={draw=main, fill=rowfill, rounded corners=2pt, align=center, font=\scriptsize,
		text width=2.15cm, execute at begin node={\hyphenpenalty=10000 \exhyphenpenalty=10000 }, minimum height=1.0cm, inner sep=2pt, text=main},
	pre/.style={draw=main, fill=infotablecolor!40, rounded corners=2pt, align=center, font=\scriptsize,
		text width=2.15cm, execute at begin node={\hyphenpenalty=10000 \exhyphenpenalty=10000 }, minimum height=1.6cm, inner sep=2pt, text=main},
	e/.style={-{Stealth[length=2.2mm]}, draw=main!80, thick},
]
\node[pre] (P) at (0,0) {Вища математика; загальна й теоретична фізика; електродинаміка};
\node[t] (T1) at (2.85,0)    {\textbf{1}. Плазма як колективне середовище};
\node[t] (T2) at (5.7,1.9)  {\textbf{2}. Колективні коливання та хвилі};
\node[t] (T3) at (5.7,0.2)  {\textbf{3}. Рух частинки в полях};
\node[t] (T4) at (5.7,-1.5) {\textbf{4}. Зіткнення, провідність, перенос};
\node[t] (T5) at (8.55,-0.65) {\textbf{5}. Гідродинаміка та МГД};
\node[t] (T6) at (8.55,1.9)  {\textbf{6}. Хвилі в замагніченій плазмі};
\node[t] (T7) at (11.4,1.9) {\textbf{7}. Кінетичний опис};
\node[t] (T8) at (11.4,0.2) {\textbf{8}. Нестійкості плазми};
\node[t] (T9) at (14.25,0.2) {\textbf{9}. Магнітне утримання};
\node[t] (T10) at (14.25,-1.5) {\textbf{10}. Керований синтез};
\node[t] (T11) at (8.55,-2.9) {\textbf{11}. Газовий розряд, низькотемп. плазма};
\node[t] (T12) at (14.25,1.9) {\textbf{12}. Масштаби, режими, межі моделей};
\draw[e] (P) -- (T1);
\draw[e] (T1.east) -- ++(0.32,0) |- (T2.west);
\draw[e] (T1.east) -- ++(0.32,0) |- (T3.west);
\draw[e] (T1.east) -- ++(0.32,0) |- (T4.west);
\draw[e] (T3.east) -- ++(0.32,0) |- (T5.west);
\draw[e] (T4.east) -- ++(0.32,0) |- (T5.west);
\draw[e] (T2) -- (T6);
\draw[e] (T5) -- (T6);
\draw[e] (T6) -- (T7);
\draw[e] (T7) -- (T8);
\draw[e] (T5) -- (T8);
\draw[e] (T8) -- (T9);
\draw[e] (T9) -- (T10);
\draw[e] (T7) -- (T12);
\draw[e] (T9) -- (T12);
\draw[e] (T4.south) |- (T11.west);
\end{tikzpicture}
\caption{Залежності між темами курсу}\label{fig:topics-graph}
\end{figure}


\section{Зміст навчальної дисципліни}

%Надається перелік розділів і тем всієї дисципліни.

Курс складається з дванадцяти розділів: плазма як іонізований газ і колективне середовище; колективні коливання та хвилі в плазмі; рух зарядженої частинки в електричному та магнітному полях; зіткнення, провідність та перенос у плазмі; гідродинамічний та магнітогідродинамічний опис; хвилі в замагніченій плазмі; кінетичний опис плазми; загальний принцип нестійкості плазми; магнітне утримання плазми; керований термоядерний синтез; низькотемпературна плазма та газовий розряд; основні масштаби, режими та межі моделей.

%Зазначається: базова (підручники, навчальні посібники) та додаткова (монографії, статті, документи, електронні ресурси) література, яку потрібно прочитати або використовувати для опанування дисципліни.
%
%Можна надати рекомендації та роз’яснення:
%\begin{itemize}
%\item де можна знайти зазначені матеріали (бібліотека, методичний кабінет, інтернет тощо);
%\item що з цього є обов’язковим для прочитання, а що факультативним;
%\item як саме студент/аспірант має використовувати ці матеріали (читати повністю, ознайомитись тощо);
%\end{itemize}
% зв’язок цих ресурсів з конкретними темами дисципліни.
%Бажано зазначати не більше п’яти базових джерел, які є вільно доступними, та не більше 20
%додаткових.


\section{Навчальні матеріали та ресурси}\label{sec:refsec}

Нижче наводиться перелік навчальних матеріалів та ресурсів для засвоєння матеріалу, розглядуваного на лекційних заняттях та для додаткового вивчення.

\nocite{%
	ponomarenko2026plasma,
	anisimov2018,
	sitenko1996,
	zagorodny2006kinetic,
	chen2016introduction,
	gurnett2017introduction,
	bellan2006fundamentals,
	freidberg2007plasma,
	stix1992waves,
	krall1973principles,
	kulsrud2005plasma,
	hazeltine2004framework,
	priest2000reconnection,
	goldston1995introduction,
	boyd2003physics,
	bittencourt2004fundamentals,
	nicholson1983introduction,
	lieberman2005principles,
	raizer1991gas,
	chen2003lecture,
}

\printbibliography[category=Main, heading=subbibliography, title={Підручники та посібники}]
\addtocategory{Main}{%
	ponomarenko2026plasma,
	anisimov2018,
	sitenko1996,
	zagorodny2006kinetic,
}

\printbibliography[category=Additional, heading=subbibliography, title={Додаткова}]
\addtocategory{Additional}{%
	chen2016introduction,
	gurnett2017introduction,
	bellan2006fundamentals,
	freidberg2007plasma,
	stix1992waves,
	krall1973principles,
	kulsrud2005plasma,
	hazeltine2004framework,
	priest2000reconnection,
	goldston1995introduction,
	boyd2003physics,
	bittencourt2004fundamentals,
	nicholson1983introduction,
	lieberman2005principles,
	raizer1991gas,
	chen2003lecture,
}

%\printbibliography[category=Problems, heading=subbibliography, title={Задачники}]
%\addtocategory{Problems}{%
%
%}


\syllabuschapter{Навчальний контент}

\section{Методика опанування навчальної дисципліни (освітнього компонента)}

\emergencystretch=3em
\textbf{Методи навчання.} Залежно від джерела подачі і сприйняття навчальної інформації використовуються словесні, наочні та практичні методи; за характером пізнавальної діяльності --- пояснювально-ілюстративний, репродуктивний, частково-пошуковий, дослідницький методи та метод проблемного викладу; за системним підходом до навчального процесу виділяють прийоми і методи організації і проведення навчання, стимулювання і мотивації здобувачів вищої освіти, контролю і самоконтролю в навчанні, а також інтегровані методи.

Інтерактивні методи навчання --- словесна і практична взаємодія викладача і здобувача. Активні методи навчання --- моделювання професійної діяльності (виконання будь-яких завдань в аудиторних умовах або виробнича практика). Інноваційні методи навчання --- оновлені традиційні методи викладання, заняття на основі інформаційно-комунікаційних технологій. Наочний метод навчання заснований на візуальному сприйнятті інформації (відеоролики, таблиці, графіки тощо). Практичний метод навчання --- прийоми і форми навчання, засновані на виконанні практичних завдань. Методи дистанційного навчання --- процес навчання на відстані від викладача за допомогою інформаційних ресурсів (онлайн-ресурси); ці методи нічим не відрізняються від традиційних, окрім онлайн-формату.

Лекційні заняття здійснюються з використанням комплексу методів, зокрема проблемного, частково-пошукового, словесного, наочного та інших. Під час проведення практичних занять і при самостійній роботі застосовуються практичний, дослідницький, наочний, словесний, проблемний, частково-пошуковий та інші групи методів. Спеціальні методи навчання: презентації, дискусія, аналітична доповідь, завдання аналітичного характеру, кейсові завдання. За джерелами набуття знань використовується поєднання словесних (монолог), наочних (демонстрація слайдів) і практичних методів навчання.

%Надається інформація (за розділами, темами) про всі навчальні заняття (лекції, практичні,
%семінарські, лабораторні) та надаються рекомендації щодо їх засвоєння (наприклад, у формі
%календарного плану чи деталізованого опису кожного заняття та запланованої роботи).

\subsection*{Лекційні та практичні заняття}

Для поглибленого вивчення питань доцільно використовувати не лише рекомендований тематичний лекційний матеріал, а і ті джерела, що вказані в розділі \ref{sec:refsec}.


% Години на тему (по одному числу на розділ, у порядку рядків таблиці); суми рахує LaTeX
% «Де вивчати» у комірці теми: тонка лінія й дрібний шрифт
% Другий (необов'язковий) аргумент --- ім'я сторінки розділу в HTML-версії книги: після «Де вивчати» клікабельне посилання (QR-коду немає: у PDF є гіперпосилання)
\newcommand{\onlineURL}{https://sergiokapone.github.io/plasma-physics/online}
\newcommand{\wherestudy}[2][]{\par\smallskip\noindent\textcolor{infotablecolor}{\rule{\linewidth}{0.4pt}}\par\smallskip
	\ifx\relax#1\relax
		{\footnotesize\textit{Де вивчати:} #2}%
	\else
		{\footnotesize\textit{Де вивчати:} #2 \href{\onlineURL/#1.html}{\textit{Читати онлайн (HTML-версія розділу).}}}%
	\fi}
\ExplSyntaxOn
\clist_const:Nn \c_syl_lec_clist  { 4,2,4,2,4,2,4,2,2,2,2,0,0 }
\clist_const:Nn \c_syl_prac_clist { 2,2,2,2,2,2,0,0,2,0,0,0,2 }
\clist_const:Nn \c_syl_srs_clist { 22,16,14,8,4,4,6 }
\NewExpandableDocumentCommand \SumSrs {} { \int_eval:n { \clist_use:Nn \c_syl_srs_clist { + } } }
\NewExpandableDocumentCommand \SrsH { m } { \clist_item:Nn \c_syl_srs_clist {#1} }
% Число годин без «--» для нуля (для рисунка з розподілом годин)
\NewExpandableDocumentCommand \LecN  { m } { \clist_item:Nn \c_syl_lec_clist  {#1} }
\NewExpandableDocumentCommand \PracN { m } { \clist_item:Nn \c_syl_prac_clist {#1} }
\NewExpandableDocumentCommand \SumLec  {}{ \int_eval:n { \clist_use:Nn \c_syl_lec_clist  { + } } }
\NewExpandableDocumentCommand \SumPrac {} { \int_eval:n { \clist_use:Nn \c_syl_prac_clist { + } } }
\NewExpandableDocumentCommand \LecH  { m } { \int_compare:nNnTF { \clist_item:Nn \c_syl_lec_clist  {#1} } = { 0 } { -- } { \clist_item:Nn \c_syl_lec_clist  {#1} } }
\NewExpandableDocumentCommand \PracH { m } { \int_compare:nNnTF { \clist_item:Nn \c_syl_prac_clist {#1} } = { 0 } { -- } { \clist_item:Nn \c_syl_prac_clist {#1} } }
\ExplSyntaxOff
\DefTblrTemplate{contfoot-text}{default}{}
\DefTblrTemplate{conthead-text}{default}{}
\DefTblrTemplate{caption}{default}{}
\DefTblrTemplate{conthead}{default}{}
\DefTblrTemplate{capcont}{default}{}
\begin{longtblr}[]{
	colspec = {Q[l, m, fg=main]X[j,m,fg=main]Q[c,m,fg=main]Q[c,m,fg=main]Q[c,m,fg=main]},
    rowhead=1,
    row{1} = {bg=rowfill},
    hlines = {1-Z}{infotablecolor},
    vlines = {1-Z}{infotablecolor},
	}
    % --------------------------------------------------------------------------
	№           &  \SetCell{c} Назва теми лекції та перелік основних питань
    & {Лекції,\\ год} & {Практичні,\\ год} & {Самостійне \\ опрацювання}                                                                   \\
    % --------------------------------------------------------------------------
    % --------------------------------------------------------------------------
	\rownumber. & \textbf{Плазма: від іонізованого газу до колективного середовища}.
    Плазма довкола нас: Всесвіт, природна й технічна плазма (оглядово). Що таке плазма, її види й способи отримання. Іонізація, формула Саха. Температура плазми, електронна та іонна компоненти. Квазінейтральність і плазмові коливання. Дебаївське екранування, дебаївська сфера, умова колективності. Слабкозв'язана та сильнозв'язана плазма. Довжина Томаса--Фермі. Шар біля стінки, критерій Бома, зонд Ленгмюра. \wherestudy[basics]{\cite[розд.~0 (огляд, §~0.1 «Що таке плазма?»), с.~14--29; розд.~1, §§~1.1--1.8, с.~30--53]{ponomarenko2026plasma}; \cite[§§~1.2--1.6, с.~2--12; §~8.2 (шар, зонд), с.~269--276]{chen2016introduction}; \cite[§§~2.1--2.2, с.~24--34; §~3.1 (Саха), с.~52--53; §~8.2 (зонди), с.~209--214]{anisimov2018}.}
    & \LecH{1} & \PracH{1} & -- \\
	\rownumber. & \textbf{Колективні коливання та хвилі в плазмі}.
    Діелектрична проникність холодної плазми. Закон дисперсії, частота відсікання. Ленгмюрівські хвилі. Зіткнення й поглинання електромагнітних хвиль. Відбиття від плазми, критична концентрація. Іоносфера і далекий радіозв'язок. Інтерферометрія і рефлектометрія. \wherestudy[waves]{\cite[розд.~2, §§~2.1--2.7, с.~57--68]{ponomarenko2026plasma}; \cite[§§~4.3--4.4, с.~78--88; §§~4.12--4.13, с.~106--113]{chen2016introduction}; \cite[§~5.1, с.~114--126; §~7.2 (відбиття), с.~171--178]{anisimov2018}.}
    & \LecH{2} & \PracH{2} & -- \\
	\rownumber. & \textbf{Рух зарядженої частинки в електричному та магнітному полях}.
    Рух в однорідному магнітному полі. Дрейфи в схрещених полях: електричний, гравітаційний, градієнтний, відцентровий. Поляризаційний дрейф. Магнітний момент і адіабатичний інваріант. Магнітне дзеркало і конус втрат. Подовжній інваріант. Диполь Землі та радіаційні пояси. Аналізатор енергій і ECE. \wherestudy[motion]{\cite[розд.~3, §§~3.1--3.8, с.~72--85]{ponomarenko2026plasma}; \cite[§§~2.2--2.8, с.~19--49]{chen2016introduction}; \cite[§§~4.1--4.2, с.~81--93]{anisimov2018}.}
    & \LecH{3} & \PracH{3} & -- \\
	\rownumber. & \textbf{Зіткнення, провідність та перенос у плазмі}.
    Кулонівські зіткнення. Провідність плазми. Теплопровідність плазми. Дифузія в плазмі, амбіполярна дифузія. Температура з опору і з ширини ліній. \wherestudy[collisions]{\cite[розд.~4, §§~4.1--4.5, с.~89--96]{ponomarenko2026plasma}; \cite[§~5.2.1, с.~148--150; §~5.5, с.~157--164; §~5.6, с.~164--171]{chen2016introduction}; \cite[§§~3.2--3.5, с.~53--80]{anisimov2018}.}
    & \LecH{4} & \PracH{4} & -- \\
	\rownumber. & \textbf{Гідродинамічний та магнітогідродинамічний опис}.
    Умови застосовності гідродинамічного опису. Дворідинна модель. Перехід до однорідинної МГД. Основне рівняння магнітної гідродинаміки. Альвенівська швидкість і число Альвена. МГД-рівновага та магнітні поверхні. Енергія електромагнітного поля та плазми. Діамагнетизм плазми. Вмороженість магнітного поля та його дифузія. Пересполучення силових ліній. \wherestudy[fluid]{\cite[розд.~5, §§~5.1--5.12, с.~100--123]{ponomarenko2026plasma}; \cite[гл.~3, с.~51--74; §§~5.7--5.9, с.~172--178; §§~6.2--6.4, с.~188--196; §~8.9 (пересполучення), с.~322--324]{chen2016introduction}; \cite[§~2.4, с.~44--51]{anisimov2018}.}
    & \LecH{5} & \PracH{5} & -- \\
	\rownumber. & \textbf{Хвилі в замагніченій плазмі}.
    Відновлювальні сили в плазмі. Звукова та іонно-звукова хвилі. Лінеаризована МГД і загальне дисперсійне рівняння. Альвенівська хвиля. Магнітозвукові хвилі. Високочастотні хвилі в замагніченій плазмі. Межі застосовності МГД-хвиль. \wherestudy[linearwaves]{\cite[розд.~6, §§~6.1--6.8, с.~127--141]{ponomarenko2026plasma}; \cite[§§~4.5--4.6, с.~89--92; §§~4.14--4.21, с.~113--144]{chen2016introduction}; \cite[§§~5.2--5.3, с.~127--141]{anisimov2018}.}
    & \LecH{6} & \PracH{6} & -- \\
	\rownumber. & \textbf{Кінетичний опис плазми}.
    Функція розподілу та фазовий простір. Моменти функції розподілу. Кінетичне рівняння, рівняння Власова. Максвеллівський розподіл і температура. Від кінетики до гідродинаміки. Лінеаризація кінетичного рівняння та хвильовий відгук. Загасання Ландау. Коли потрібен кінетичний опис. Томсонівське розсіяння. \wherestudy[kinetics]{\cite[розд.~7, §§~7.1--7.10, с.~144--161]{ponomarenko2026plasma}; \cite[§§~7.1--7.6, с.~211--244]{chen2016introduction}; \cite[§~2.3, с.~35--43; §~8.4 (розсіяння), с.~225]{anisimov2018}.}
    & \LecH{7} & \PracH{7} & + \\
	\rownumber. & \textbf{Загальний принцип нестійкості плазми}.
    Лінеаризація й дисперсійне рівняння. Джерела енергії нестійкості. Рівновага та стійкість у магнітному полі. Лінійна та нелінійна стадії. Нестійкість Релея--Тейлора та жолобкова нестійкість. Двопотокова нестійкість і нестійкість пучка. Нестійкості анізотропної плазми. Енергетичний принцип. Дрейфові хвилі та аномальний перенос. \wherestudy[instability]{\cite[розд.~8, §§~8.1--8.15, с.~165--178]{ponomarenko2026plasma}; \cite[§§~6.5--6.9, с.~196--210]{chen2016introduction}; \cite[§§~4.3--4.5, с.~94--105; §§~7.5--7.6, с.~188--201]{anisimov2018}.}
    & \LecH{8} & \PracH{8} & + \\
	\rownumber. & \textbf{Магнітне утримання плазми}.
    Проблема утримання. Відкриті пастки, магнітне дзеркало. Пінч-ефект, рівновага $z$-пінча, співвідношення Беннета. Рівновага $\theta$-пінча. Нестійкості пінча. Тороїдальні системи: токамак і стеларатор. Захоплені частинки й неокласичний перенос. Зсув Шафранова та межа за тиском. \wherestudy[confinement]{\cite[розд.~9, §§~9.1--9.10, с.~182--192]{ponomarenko2026plasma}; \cite[§~10.2.1, с.~356--386; §~5.10, с.~178--185; §~6.3, с.~190--192]{chen2016introduction}; \cite[§§~4.2, 4.4--4.5, с.~85--93 і~96--105]{anisimov2018}.}
    & \LecH{9} & \PracH{9} & -- \\
	\rownumber. & \textbf{Керований термоядерний синтез}.
    Реакції та паливо. Кулонівський бар'єр, тунелювання, фактор Гамова. Баланс потужностей і критерій Лоусона. Втрати. Магнітне й інерційне утримання. Від реакції до реактора. Інерційний термоядерний синтез. Сучасний стан досліджень. Нейтронна діагностика. \wherestudy[fusion]{\cite[розд.~10, §§~10.1--10.10, с.~195--204]{ponomarenko2026plasma}; \cite[§~1.7.2, с.~13--14; §~10.2, с.~356; §~10.4, с.~393--399]{chen2016introduction}; \cite[§~4.6, с.~106--113]{anisimov2018}.}
    & \LecH{10} & \PracH{10} & + \\
    % --------------------------------------------------------------------------
	\rownumber. & \textbf{Низькотемпературна плазма та газовий розряд}.
    Іонізація та її швидкість. Пробій газу: механізм Таунсенда і закон Пашена. Тліючий розряд і дуга. Приелектродний шар у розряді. Глобальна модель розряду. Високочастотні джерела плазми. Плазмова обробка матеріалів. Оптична емісія і ВЧ-зонди. \wherestudy[discharge]{\cite[розд.~11, §§~11.1--11.8, с.~207--218]{ponomarenko2026plasma}; \cite[§~1.7.1, с.~13; §~8.2, с.~269--276; §~10.5, с.~399--404; §~9.5, с.~346--349]{chen2016introduction}; \cite[§~1.2, с.~19--23; §~8.2, с.~209--214]{anisimov2018}.}
    & \LecH{11} & \PracH{11} & + \\
	\rownumber. & \textbf{Основні масштаби, режими та межі моделей}.
    Просторові та часові масштаби. Зіткнювальний і беззіткнювальний режими. Вибір моделі за масштабами. Ієрархія масштабів у магнітному утриманні. Від кінетики до МГД. Безрозмірні параметри МГД. Межі простих моделей. Що вимірюють на якому масштабі. \wherestudy[scales]{\cite[розд.~12, §§~12.1--12.9, с.~221--229]{ponomarenko2026plasma}; \cite[§§~1.5--1.6, с.~11--12; §~3.6, с.~73--74; §~4.7, с.~92--93]{chen2016introduction}; \cite[§~3.3, с.~57--62; §~8.1, с.~202--208]{anisimov2018}.}
    & \LecH{12} & \PracH{12} & + \\
    % --------------------------------------------------------------------------
	& \textbf{Модульна контрольна робота}. Письмова робота на практичному занятті за матеріалом лекцій і практичних занять.
    & \LecH{13} & \PracH{13} & -- \\
    % --------------------------------------------------------------------------
	\SetCell[c=2]{r} \textbf{Разом} & & \textbf{\SumLec} & \textbf{\SumPrac} & \textbf{--} \\
    % --------------------------------------------------------------------------
\end{longtblr}

\Needspace{6\baselineskip}
Послідовність і обсяг аудиторної роботи показано на рис.~\ref{fig:hours-strip}: ширина кожного блока пропорційна кількості годин, а номер над ним --- номер теми в таблиці вище.

\begin{figure}[h!]
\centering
\begin{tikzpicture}[x=1cm, y=1cm, font=\scriptsize]
\newcommand{\advx}[1]{\pgfmathparse{\xx+#1}\xdef\xx{\pgfmathresult}}% \foreach працює в групі: потрібен глобальний \xx
\def\hu{0.352}% см на одну годину: 46 год ≈ 16,2 см
\gdef\xx{0}
\foreach \i in {1,...,13}{%
	\pgfmathtruncatemacro{\hl}{\LecN{\i}}%
	\pgfmathtruncatemacro{\hp}{\PracN{\i}}%
	\pgfmathsetmacro{\xs}{\xx}%
	\ifnum\hl>0
		\fill[maintext!75] (\xx,0) rectangle +(\hl*\hu,0.55);
		\node[white] at (\xx+\hl*\hu/2,0.275) {\hl};
		\advx{\hl*\hu}%
	\fi
	\ifnum\hp>0
		\ifnum\i=13 \def\pc{airforceblue!35}\else\def\pc{infotablecolor!55}\fi
		\fill[\pc] (\xx,0) rectangle +(\hp*\hu,0.55);
		\node[main] at (\xx+\hp*\hu/2,0.275) {\hp};
		\advx{\hp*\hu}%
	\fi
	\pgfmathtruncatemacro{\hsum}{\hl+\hp}%
	\ifnum\hsum>0
	\draw[main!60, thin] (\xs,0) -- (\xs,0.75);
	\ifnum\i=13
		\node[main, anchor=south] at (\xs+\hp*\hu/2,0.78) {МКР};
	\else
		\node[main, anchor=south] at (\xs+\hl*\hu/2+\hp*\hu/2,0.78) {\i};
	\fi
	\fi
}
\draw[main!60, thin] (\xx,0) -- (\xx,0.75);
\draw[main, thick] (0,0) rectangle (\xx,0.55);
\draw[-{Stealth[length=2mm]}, main, thick] (\xx,0.275) -- ++(0.55,0) node[right] {залік};
% легенда
\fill[maintext!75] (0,-0.75) rectangle +(0.35,0.3); \node[anchor=west] at (0.4,-0.6) {лекції (\SumLec\ год)};
\fill[infotablecolor!55] (4,-0.75) rectangle +(0.35,0.3); \node[anchor=west] at (4.4,-0.6) {практичні заняття (\SumPrac\ год)};
\fill[airforceblue!35] (9,-0.75) rectangle +(0.35,0.3); \node[anchor=west] at (9.4,-0.6) {модульна контрольна робота};
\end{tikzpicture}
\caption{Розподіл аудиторних годин за темами курсу}\label{fig:hours-strip}
\end{figure}


\section{Самостійна робота здобувача вищої освіти}

Самостійна робота здобувачів вищої освіти має на меті розвиток творчих здібностей та активізація їх розумової діяльності, формування потреби безперервного самостійного поповнення знань та розвиток морально-вольових зусиль. Завданням самостійної роботи здобувачів вищої освіти є навчити здобувачів вищої освіти самостійно працювати з літературою, творчо сприймати навчальний матеріал і осмислювати його та формування навичок до щоденної роботи з метою одержання та узагальнення знань, умінь і навичок.


На самостійну роботу відводяться наступні види завдань (орієнтовний розподіл годин):

\setcounter{magicrownumbers}{0}
\begin{longtblr}{
	colspec = {Q[c, m, fg=main]X[j,m,fg=main]Q[c, m, fg=main]},
	rowhead=1,
	row{1} = {bg=rowfill},
	hlines = {1-Z}{infotablecolor},
	vlines = {1-Z}{infotablecolor},
}
	№ & \SetCell{c} Вид самостійної роботи & Години \\
	\rownumber. & Самостійне опрацювання розділів, що не розбираються на практичних заняттях: загальний принцип кінетичний опис (розд.~7), нестійкості (розд.~8), термоядерний синтез (розд.~10), низькотемпературна плазма й газовий розряд (розд.~11) і масштаби та межі моделей (розд.~12): читання за посібником і за літературою, вказаною в таблиці тем, конспект ключових формул і оцінок. & \SrsH{1} \\
	\rownumber. & Розв’язування задач із розділу «Задачі» посібника до кожного практичного заняття (розд.~1--6 і~9); відповіді наведено в розділі «Відповіді» кінця посібника, свій розв’язок слід звіряти з ними й виправляти. & \SrsH{2} \\
	\rownumber. & Домашні завдання: розрахункові задачі й числові оцінки (дебаївський радіус, параметр плазми, дрейфи, критерій Лоусона тощо) за матеріалом розд.~1--10. & \SrsH{3} \\
	\rownumber. & Відповіді на контрольні питання наприкінці кожного розділу посібника для самоперевірки після лекції. & \SrsH{4} \\
	\rownumber. & Робота з інтерактивними симуляціями до розд.~1--3 (дебаївське екранування, зонд Ленгмюра, плазмові коливання, магнітне дзеркало), які розміщено в онлайн-версії посібника; параметри змінюють самостійно й порівнюють результат з оцінками з тексту. & \SrsH{5} \\
		\rownumber. & Підготовка до модульної контрольної роботи: повторення означень, виведень, оцінок і діапазонів застосовності за матеріалом лекцій і практичних занять. & \SrsH{6} \\
	\rownumber. & Підготовка до складання заліку: повторення всього курсу, розв’язування задач і відповіді на контрольні питання розділів посібника. & \SrsH{7} \\
	\SetCell[c=2]{r} \textbf{Разом} & & \textbf{\SumSrs} \\
\end{longtblr}




\Needspace{12\baselineskip}
\syllabuschapter{Політика та контроль}

\section{Політика навчальної дисципліни (освітнього компонента)}

\subsection*{Відвідування занять}
Відвідування лекцій, а також відсутність на них, не оцінюється. Однак, здобувачам вищої освіти бажано відвідувати заняття, оскільки на них викладається теоретичний матеріал та розвиваються навички, необхідні для успішного складання \control у.

\subsection*{Пропущені контрольні заходи}

Результат модульної контрольної роботи для здобувача вищої освіти, який не з’явився на контрольний захід, є нульовим. У такому разі, здобувач вищої освіти має можливість написати модульну контрольну роботу, але максимальний бал за неї буде дорівнювати 50~\% від загальної кількості балів. Повторне написання модульної контрольної роботи не допускається.

\subsection*{Календарний рубіжний контроль}

Календарний контроль проводиться з метою моніторингу виконання здобувачами індивідуальних навчальних планів згідно з графіком навчального процесу. Умовою позитивної атестації є значення поточного рейтингу здобувача вищої освіти не менше 50\% від максимально можливого на час атестації. Бал, необхідний для отримання позитивного календарного контролю, доноситься до здобувачів вищої освіти викладачем не пізніше ніж за 2 тижні до початку календарного контролю.

\subsection*{Визнання  результатів  навчання,  набутих  у  неформальній/інформальній освіті}

Порядок визнання таких результатів регламентується Положенням про визнання результатів навчання, набутих у неформальній / інформальній освіті (\url{https://osvita.kpi.ua/index.php/node/179}).  Можуть бути зараховані окремі змістовні модулі або теми дисципліни. В такому разі здобувач звільняється від виконання відповідних завдань, отримуючи за них максимальний бал відповідно до рейтингової системи оцінювання.


\subsection*{Політика щодо академічної доброчесності}

Обов’язковою умовою виконання завдань з освітньої компоненти є дотримання політики та принципів академічної доброчесності (\url{https://kpi.ua/academic-integrity}), які, у тому числі, викладено у Кодексі честі Національного технічного університету України «Київський політехнічний інститут імені Ігоря Сікорського» (\url{https://kpi.ua/code}), Положенні про систему запобігання академічному плагіату в КПІ ім. Ігоря Сікорського (\url{https://osvita.kpi.ua/node/47}). У разі виявлення дублювання робіт, плагіату роботи здобувачі отримують нульовий рейтинг.

\subsection*{Політика використання штучного інтелекту}

Використання штучного інтелекту (далі, ШІ) регламентується «Політикою використання штучного інтелекту для академічної діяльності в КПІ ім. Ігоря Сікорського» (\url{https://osvita.kpi.ua/node/1225}). Усі навчальні завдання з дисципліни мають бути результатом власної оригінальної роботи здобувача.

\subsection*{Норми етичної поведінки}

Норми етичної поведінки здобувачів вищої освіти і працівників визначені у розділі 2 Кодексу честі Національного технічного університету України «Київський політехнічний інститут імені Ігоря Сікорського». Детальніше: \url{https://kpi.ua/code}.

\subsection*{Процедура оскарження результатів контрольних заходів}

Здобувачі вищої освіти мають можливість підняти будь-яке питання, яке стосується процедури контрольних заходів та очікувати, що воно буде розглянуто згідно із наперед визначеними процедурами (згідно <<Положення про систему забезпечення якості вищої освіти у Національному технічному університеті України «Київський політехнічний інститут імені Ігоря Сікорського>>, <<Положення про організацію навчального процесу>>).


\section{Рейтингова система оцінювання результатів навчання (РСО)}

\subsection*{Форма семестрового контролю: залік}

\textbf{Тип РСО: РСО-1 (відповідно до п. 2.3, 3.8–3.14 \href{https://osvita.kpi.ua/sites/default/files/downloads/Pologennia_RSO_2025.pdf}{Положення про систему оцінювання результатів навчання КПІ, 2025)}.}

\subsection*{Загальні принципи оцінювання}

Результати навчання з дисципліни (освітнього компонента) оцінюються за рейтинговою системою оцінювання першого типу (РСО-1), семестровий контроль з якої передбачений у формі заліку, і включає оцінювання заходів поточного контролю з дисципліни впродовж семестру, а рейтингова оцінка здобувача складається з балів, отриманих здобувачем за результатами заходів поточного контролю. Оцінювання результатів навчання здійснюється за 100-бальною шкалою.

Здобувачі, які набрали менше 60 балів, а також ті, які бажають підвищити свою підсумкову оцінку, на заліковому тижні проходять семестровий контроль у вигляді залікової контрольної роботи. Під час складання заліку застосовується \textbf{м’який підхід}: якщо результат заліку нижчий, ніж підсумок поточного контролю, до відомості вноситься вищий результат.

\subsection*{Види контролю та їхня вага у рейтинговій системі}

\begin{table}[h!]
\centering
\caption{Розподіл балів за видами контролю з дисципліни «\discipline» (РСО-1, залік)}
\begin{tblr}{
colspec = {Q[l,m, 7cm]|Q[c]|Q[c]|},
row{1}={c,m, fg=main, bg=rowfill},
        hlines = {1-Z}{infotablecolor},
        vlines = {1-Z}{infotablecolor},
}
{Вид контролю} & {Форма проведення} & {Максимальна\\ кількість балів} \\
Модульна контрольна робота (перевірка теоретичного матеріалу, розрахункові задачі) & Синхронна & 50 \\
Домашні завдання & Асинхронна & 40 \\
Поточна активність, короткі тести, участь у семінарах & Синхронна / асинхронна & 10 \\
\SetCell[r=1,c=2]{r}\textbf{Загальна кількість балів (поточний контроль)} &  & \textbf{100} \\
\end{tblr}
\end{table}

\subsection*{Підсумкове оцінювання}

\noindent Якщо здобувач вищої освіти набрав сумарно $R_{\text{поточн.}} \ge 60$ балів, він отримує позитивну рейтингову оцінку без проходження семестрового контролю.

У разі набрання $R_{\text{поточн.}} < 60$ балів, здобувач вищої освіти складає залік у формі залікової контрольної роботи, яка оцінюється за 100-бальною шкалою.
Підсумковий рейтинг визначається за формулою:
\begin{equation*}
R_\text{підсумк.} =
\begin{cases}
R_\text{поточн.}, & \text{якщо } R_\text{поточн.} \ge 60; \\[6pt]
\max\bigl(R_\text{поточн.},\,R_\text{залік}\bigr), & \text{якщо } R_\text{поточн.} < 60.
\end{cases}
\end{equation*}

\subsection*{Питання до залікової контрольної роботи}

Залікова контрольна робота складається з теоретичних питань і розрахункових задач. Теоретичні питання відповідають розділам «Контрольні питання» наприкінці розділів посібника \cite{ponomarenko2026plasma}, розрахункові задачі --- розділам «Задачі» тих самих розділів (відповіді наведено в кінці посібника). Сторінки «Контрольних питань»: 
% BEGIN QUESTIONS (генерується update_refs.py)
розд.~0 --- с.~29; розд.~1 --- с.~55; розд.~2 --- с.~70; розд.~3 --- с.~88; розд.~4 --- с.~98; розд.~5 --- с.~126; розд.~6 --- с.~143; розд.~7 --- с.~163; розд.~8 --- с.~180; розд.~9 --- с.~193; розд.~10 --- с.~206; розд.~11 --- с.~219; розд.~12 --- с.~231.
% END QUESTIONS

Склад рейтингу та порядок отримання підсумкової оцінки показано на рис.~\ref{fig:rso-scheme}.

\begin{figure}[h!]
\centering
\begin{tikzpicture}[x=1cm, y=1cm, font=\footnotesize,
	b/.style={draw=main, fill=rowfill, rounded corners=2pt, align=center, inner sep=3pt, text=main, minimum height=0.9cm},
	d/.style={draw=main, fill=infotablecolor!40, diamond, aspect=2.6, align=center, inner sep=1pt, text=main},
	e/.style={-{Stealth[length=2.2mm]}, draw=main!80, thick}]
% склад 100 балів (шкала 0.16 см на бал)
\def\sc{0.16}
\fill[maintext!75] (0,0) rectangle (50*\sc,0.7);
\fill[infotablecolor!70] (50*\sc,0) rectangle (90*\sc,0.7);
\fill[airforceblue!35] (90*\sc,0) rectangle (100*\sc,0.7);
\draw[main, thick] (0,0) rectangle (100*\sc,0.7);
\node[white] at (25*\sc,0.35) {МКР: 50};
\node[main] at (70*\sc,0.35) {Домашні завдання: 40};
\node[main] at (95*\sc,0.35) {10};
\node[main, anchor=north, align=center] at (95*\sc,-0.05) {\scriptsize активність,\\[-2pt]\scriptsize тести};
\draw[main, dashed] (60*\sc,-0.15) -- (60*\sc,0.95);
\node[main, anchor=south] at (60*\sc,0.95) {60: поріг заліку};
\node[main, anchor=north] at (0,-0.05) {0};
\node[main, anchor=north] at (50*\sc,-0.05) {50};
% схема отримання оцінки
\node[b] (A) at (2.2,-2.2) {Поточний контроль:\\ $R_{\text{поточн.}}$};
\node[d] (Q) at (6.3,-2.2) {$R_{\text{поточн.}}\ge 60$?};
\node[b] (Y) at (11.4,-1.4) {Оцінка $=R_{\text{поточн.}}$\\ (залік без складання)};
\node[b] (N) at (11.4,-3.2) {Залікова контрольна\\ робота, 0--100 балів};
\node[b] (F) at (11.4,-4.7) {Підсумок: більший із двох\\ результатів (м'який підхід)};
\draw[e] (A) -- (Q);
\draw[e] (Q.east) -- ++(0.35,0) |- node[near end, above, font=\scriptsize] {так} (Y.west);
\draw[e] (Q.east) -- ++(0.35,0) |- node[near end, below, font=\scriptsize] {ні} (N.west);
\draw[e] (N) -- (F);
\draw[dashed, main!80, thick] (Y.east) -- ++(0.6,0) coordinate (c);
\draw[e, dashed] (c) -- node[right, font=\scriptsize, align=left] {для\\ підви-\\ щення} (c |- N.east) -- (N.east);
\end{tikzpicture}
\caption{Склад рейтингу (РСО-1) та отримання підсумкової оцінки}\label{fig:rso-scheme}
\end{figure}




%\section{Види контролю та рейтингова система оцінювання результатів навчання (РСО)}
%
%%Вказуються всі види контролю та бали за кожен елемент контролю, наприклад:
%%Поточний контроль: експрес-опитування, опитування за темою заняття, МКР, тест тощо
%%Календарний контроль: провадиться двічі на семестр як моніторинг поточного стану
%%виконання вимог силабусу.
%%
%%Семестровий контроль: екзамен / залік / захист курсового проекту (роботи)
%%Умови допуску до семестрового контролю: мінімально позитивна оцінка за індивідуальне
%%завдання / зарахування усіх лабораторних робіт / семестровий рейтинг більше ХХ балів.
%
%Видами контролю успішності засвоєння матеріалу дисципліни є  модульна контрольна робота (МКР), домашня контрольна робота та семестровий контроль.
%
%\subsection*{Практичні заняття}
%
%\pgfmathsetmacro{\pracBal}{2.5}
%\pgfmathsetmacro{\pracNum}{8}
%Практичні заняття оцінюються на основі розв'язаних на парі практичних задач. Задачы задаються в classroom. Обчислюється середній бал за всі здані практичні задачі, який множиться на ваговий коефіцієнт $0.2$. Максимальний бал за кожне практичне заняття \pracBal{} (з урахуванням вагового коефіцієнту), максимальний бал
%\xinttheexpr\pracBal*\pracNum\relax.
%%\pgfmathparse{\pracBal*\pracNum}\pgfmathprintnumber[precision=0]\pgfmathresult.
%
%Бал за задачі, які були здані пізніше вказаного в classroom терміну, знижується на 20\%.
%
%Задачі бажано оформлювати у форматі \texttt{PDF}, який створюється за допомогою \LaTeX.
%
%
%
%\subsection*{Модульна контрольна робота}
%
%\pgfmathsetmacro{\mkrBal}{10}
%
%Модульна контрольна робота проводиться після завершення першої частини курсу <<\discipline>> проводиться протягом 2-х академічних годин на практичних заняттях. Вона складається з $3$-х задач і передбачає письмовий розв’язок задачі, подібних до тих, що розглядались на практичних заняттях та під час виконання домашніх робіт.
%
%Критерії оцінювання модульної контрольної роботи:
%\begin{enumerate}[label=$\bullet$]
%	\item максимальна кількість балів за кожну задачу --- повна правильна відповідь, $95$\% інформації, там де треба наведено рисунки, позначення, є письмовий коментар щодо базових понять та законів, які використовуються під час розв’язку задачі,
%	\item $75$\%  --- розв’язок правильний, не всі умови попереднього пункту виконано,
%	\item $60$\%  --- наведено основні базові поняття для розв’язку, розв’язок неправильний.
%	\item списані відповіді, які студент не може пояснити, не зараховуються.
%\end{enumerate}


%\subsection*{Реферат}
%
%Тема роботи обирається студентом за власним бажанням і узгоджується з викладачем або обирається із списку (див. розділ~\ref{sec:referat}), який надає викладач. Тема роботи не повинна дублювати лекційний матеріал. Найбільш вітається робота, яка містить критичний аналіз оригінальних наукових публікацій. Допускається детальний аналіз тем, викладених в посібниках, але не освітлених в лекційному курсі. Оцінка визначається комплексом: вибір тематики, оформлення, творчий підхід. Фінальна оцінка складається на основі захисту --- публічної співбесіди з викладачем у присутності студентів. В обговоренні реферату можуть брати участь усі студенти групи. Реферат оцінюється за наступними критеріями Якість викладення результатів (макс. $30$ балів) Якість оформлення роботи (макс. $20$ балів). Захист реферату (макс. $50$ балів). Ваговий коефіцієнт $0.3$.
%
%\pgfmathsetmacro{\referatBal}{30}



%\subsection*{Умови допуску до \control у}
%
%В таблиці наведені умови допуску до семестрового контролю.
%
%\begin{center}\setcounter{magicrownumbers}{0}
%	\begin{tblr}{
%        colspec={|Q[c,m,fg=main]|Q[l,m,fg=main]|Q[l,m,fg=main]|},
%        row{1}={rowfill},
%        hlines = {1-Z}{infotablecolor},
%        vlines = {1-Z}{infotablecolor},
%    }
%	    №      & Обов’язкова умова допуску до \control у & Критерій  \\
%	    \rownumber & Потчний рейтинговий бал             & $\ge 30$  \\
%	    \rownumber & МКР                                 & виконана  \\
%	    \rownumber & Реферат                             & захищений \\
%	\end{tblr}%
%\end{center}

%\subsection*{Семестровий контроль (\control)}
%
%\pgfmathsetmacro{\controlBal}{40}
%
%\Control~--- письмова відповідь на \emph{\color{main} два} теоретичних питання білету.
%
%Оцінка за \control\ складається з рейтингу за семестр та оцінки за письмову відповідь напитання білету (максимум $40$ баллів).
%
%Оцінка формується на основі наступних критеріїв:
%\begin{enumerate}[label=$\bullet$]
%	\item від $35$ до $40$ --- повна правильна відповідь, $95$\% інформації, наведено рисунки, позначення, є письмовий коментар щодо базових понять та законів, формулювання та терміни точні, терміни роз’яснено, повна правильна відповідь на уточнюючі запитання
%	\item від $25$ до $35$ --- правильна відповідь, 80\% інформації, наведено рисунки, позначення, є письмові коментарі щодо базових понять та законів, формулювання та терміни по суті правильні але не повні, терміни роз’яснено, правильна відповідь на уточнюючі запитання
%	\item від $20$ до $25$ --- по суті правильна але неповна відповідь, $70$\% інформації, наведено рисунки та позначення, відсутні письмові коментарі щодо базових понять та законів, формулювання та терміни по суті правильні але не повні, терміни не роз’яснено, правильна відповідь на більшість уточнюючих запитання
%	\item від $10$ до $20$ --- відповідь неповна, $50$\% інформації, не наведено рисунки та позначення, відсутні письмові коментарі щодо базових понять та законів, формулювання та терміни в основному правильні але не повні, терміни не роз’яснено, відповіді на уточнюючі запитання не повні
%	\item від $0$ до $10$ --- відповідь неповна, $30$\% інформації, не наведено рисунки та позначення, відсутні письмові коментарі щодо базових понять та законів, формулювання та терміни в основному не повні, терміни не роз’яснено, відповіді на уточнюючі запитання не повні або відсутні
%\end{enumerate}

%
%Остаточна оцінка $\mathbf{R}$ є сумою рейтингових балів отриманих за поточний контроль та балів отриманих на \control і після співбесіди зі студентом.
%
%
%\begin{center}\setcounter{magicrownumbers}{0}
%	\begin{spreadtab}{{tblr}
%    {colspec={Q[c,m,fg=main]Q[l,m,fg=main]Q[c,m,fg=main]Q[c,m,fg=main]Q[c,m,fg=main]},
%    row{1}={rowfill},
%    hlines = {1-Z}{infotablecolor},
%    vlines = {1-Z}{infotablecolor},
%    }
%    }
%	    @ №                       & @ Контрольний захід          & @ Бал       & @ Кількість                    & @ Всього                    \\
%	    @ \rownumber              & @ Практичні заняття          & \pracBal    & \pracNum                       & \STcopy{v}{c2*d2}           \\
%	    @ \rownumber              & @ Модульна контрольна робота & \mkrBal     & 1                              &                              \\
%	    @ \rownumber              & @ ДКР                    & 30 & 1                              &                              \\
%	    @ \rownumber              & @ \Control                   & \controlBal & 1                              &                              \\
%	    @\SetCell[c=4]{h,c}Всього &                              &                                &             & sum(e1:[0,-1]) \\
%	\end{spreadtab}
%\end{center}






\begin{samepage}
Таблиця відповідності рейтингових балів оцінкам за університетською шкалою.\par
\begin{center}
	\begin{tblr}{
        colspec={Q[c,m,fg=main]Q[c,m,fg=main, 0.25\linewidth]},
        row{1}={rowfill},
        hlines = {1-Z}{infotablecolor},
        vlines = {1-Z}{infotablecolor},
    }
	            Значення рейтингу              & Оцінка ECTS  \\
	       $95 \le \mathbf{R} \le 100$         &   відмінно   \\
	         $85 \le \mathbf{R} < 95$          &  дуже добре  \\
	         $75 \le \mathbf{R} < 85$          &    добре     \\
	         $65 \le \mathbf{R} < 75$          &  задовільно  \\
	         $60 \le \mathbf{R} < 65$          &  достатньо   \\
	            $ \mathbf{R} < 60$             & незадовільно \\
	    {Не здана МКР,\\ не виконані домашні завдання} & не допущено  \\
	\end{tblr}%
\end{center}\end{samepage}

%\clearpage

%\section{Додаткова інформація з дисципліни}
%
%\subsection*{Перелік рекомендованих тем рефератів}\label{sec:referat}
%\begin{enumerate}
%\item Газові розряди: їх види, фізика процесів та перспективи застосування.
%\item Плазмові двигуни для космічних апаратів: принципи роботи, типи, перспективи.
%\item Досягнення в діагностиці плазми: методи, прилади, принципи вимірювання.
%\item Плазмові лазери: особливості генерації, посилення та керування плазмовим випромінюванням.
%\item Плазмова обробка матеріалів: методи, застосування, ефекти.
%\item Плазмова нанотехнологія: синтез, модифікація та застосування наночастинок і наноструктур з використанням плазми.
%\item Плазми в електроніці: створення ефективних сенсорів та пристроїв зберігання інформації.
%\item Використання плазми для вирішення проблем екології.
%\item Плазмові джерела електричної енергії: принципи, проблеми, сучасний стан розвитку та перспективи.
%\item Термоядерна енергетика: принципи, проблеми, перспективи.
%\item Плазмова біологія та медицина: вплив плазми на живі організми, біомедичне застосування плазми.
%\item Плазма в астрофізиці: сонячний вітер та його вплив на магнітосферу планет.
%\end{enumerate}











%\begin{itemize}
%\item перелік питань, які виносяться на семестровий контроль (наприклад, як додаток до силабусу);
%\item  можливість зарахування сертифікатів проходження дистанційних чи онлайн курсів за відповідною тематикою;
%\item  інша інформація для студентів/аспірантів щодо особливостей опанування навчальної дисципліни.
%\end{itemize}

%\section{Додаткова інформація з дисципліни <<\discipline>>}


%•	перелік питань, які виносяться на семестровий контроль (наприклад, як додаток до силабусу);
%•	можливість зарахування сертифікатів проходження дистанційних чи онлайн курсів за відповідною тематикою;
%•	інша інформація для студентів/аспірантів щодо особливостей опанування навчальної дисципліни.

%\subsection*{Методи навчання}
%
%Методом проведення лекційних занять є наочний метод, який включає демонстрацію ключових дослідів та використання ілюстративного матеріалу; проблемний виклад, який дає уявлення про набуття нових знань на основі відомих фактів; бесіди і дискусії, для встановлення діалогу з аудиторією та залучення їх для активного осмислення інформації.
%
%Методами проведення практичних занять є метод тренінгу та репродуктивний метод. Для лабораторних занять використовуються частково-пошукові та дослідницькі методи.

%По всім видам занять проводяться дистанційні консультації.

%\subsection*{Перелік питань, які виносяться на усний \control}
%
%\loaditemsfromfile[enumerate]{ElMagQ.dat}
%
\Needspace{14\baselineskip}
\phantom{text}
\vfill
\thispagestyle{empty}
\printrequisites
\end{document}



